\documentclass[../../../include/open-logic-section]{subfiles}

\begin{document}
\olfileid[pt]{sth}{ordinals}{zfm}	
\olsection{$\ZFminus$: um marco}

A questão de como justificar o Esquema de Substituição (se é que o
justificamos) não é linear. Como tal, a reservamos para o
\olref[replacement][]{chap}. No entanto, com a adição da Substituição,
atingimos outro marco importante. Dispomos agora de todos os axiomas
necessários para a teoria $\ZFminus$. Em detalhe:

\begin{defn}
A teoria $\ZFminus$ tem estes axiomas: Extensionalidade, União, Pares,
Conjuntos das Partes, Infinito, e todas as instâncias dos esquemas de
Separação e Substituição. Por outras palavras, $\ZFminus$ acrescenta a
Substituição a $\Zminus$.
\end{defn}
\noindent
Isto designa a teoria de conjuntos de \emph{Zermelo--Fraenkel}
(\emph{menos} algo que veremos mais adiante). Fraenkel recebe a honra,
uma vez que lhe é atribuída a formulação da Substituição em
\citeyear{Fraenkel1922}, embora a primeira formulação precisa seja devida
a \citet{Skolem1922}.

\end{document}
